\appendix
\chapter{Pseudocódigos auxiliares}
\label{anexo:Pseudocodigo}

\section{Pseudocódigo del procesamiento de imágenes}
\label{apend:procesa}
\begin{center}
\captionof{algorithm}{Pseudo código del procesamiento de imágenes}
\label{alg:pdi}
\end{center}
\vspace{-1.6cm}
\rule{\textwidth}{0.4pt}
\vspace{-0.9cm}

\begin{algorithmic}[1]

\Statex \textbf{ENTRADA:}
\Statex fotograma\_crudo
\Comment{imagen capturada directamente por la cámara}

\Statex
\Statex \textbf{SALIDA:}
\Statex imagen\_lista
\Comment{imagen normalizada y preprocesada}

\Statex
\Statex \textbf{PROCESO}
\State Verificar que el fotograma no esté vacío

\If{fotograma\_crudo está vacío}

    \State registrar error y terminar

\EndIf

\Statex
\State Redimensionar la imagen al tamaño objetivo (ancho x alto)

\State imagen\_redimensionada $\gets$
redimensionar(fotograma\_crudo, ancho, alto)

\Statex
\State Convertir a espacio de color normalizado (escala [0, 1])

\State imagen\_normalizada $\gets$
imagen\_redimensionada / 255.0

\Statex
\State Aplicar igualación de histograma para mejorar contraste

\State imagen\_ecualizada $\gets$
ecualizar\_histograma(imagen\_normalizada)

\Statex
\State Aplicar filtro gaussiano para suavizar ruido

\State imagen\_lista $\gets$
filtro\_gaussiano(imagen\_ecualizada, sigma=0.5)

\Statex
\State \Return imagen\_lista

\end{algorithmic}
\vspace{-0.9cm}
\rule{\textwidth}{0.4pt}

\section{Pseudocódigo del módulo de Ray Casting}
\label{apend:ray}
\begin{center}
\captionof{algorithm}{Pseudocódigo del módulo de Ray Casting}
\label{alg:raycasting}
\end{center}
\vspace{-1.6cm}
\rule{\textwidth}{0.4pt}
\vspace{-0.9cm}

\begin{algorithmic}[1]

\Statex \textbf{ENTRADA:}

\Statex pose
\Comment{[x, y, theta] del usuario}

\Statex angulos\_rayos
\Comment{lista de ángulos relativos a theta}

\Statex distancias
\Comment{distancia de impacto para cada rayo}

\Statex \hspace{0.5cm}
// Variables a inicializar

\Statex reticula
\Comment{mapa de ocupación en log-odds}

\Statex resolucion
\Comment{tamaño de cada celda en metros}

\Statex incremento\_libre
\Comment{incremento log-odds para celda libre (negativo)}

\Statex incremento\_ocup
\Comment{incremento log-odds para celda ocupada (positivo)}

\Statex lo\_min, lo\_max
\Comment{límites de saturación}

\Statex
\Statex \textbf{SALIDA:}

\State reticula
\Comment{mapa actualizado}

\Statex
\Statex \textbf{PROCESO}
\For{cada rayo i EN angulos\_rayos}

    \State Calcular dirección absoluta del rayo

    \State angulo\_abs $\gets$
    pose.theta + angulos\_rayos[i]

    \State distancia\_impacto $\gets$
    distancias[i]

    \Statex

    \State Obtener secuencia de celdas con algoritmo de Bresenham

    \State celda\_origen $\gets$
    convertir\_a\_celda(pose.x, pose.y, resolucion)

    \State celda\_fin $\gets$
    convertir\_a\_celda(

    \Statex \hspace{1.3cm}
    pose.x + distancia\_impacto * coseno(angulo\_abs),

    \Statex \hspace{1.3cm}
    pose.y + distancia\_impacto * seno(angulo\_abs),

    \Statex \hspace{1.3cm}
    resolucion)

    \State celdas\_rayo $\gets$
    bresenham(celda\_origen, celda\_fin)

    \Statex

    \State Marcar todas las celdas antes del impacto como libres

    \For{cada celda c EN celdas\_rayo excepto la última}

        \State reticula[c] $\gets$
        saturar(reticula[c] + incremento\_libre,

        \Statex \hspace{2.2cm}
        lo\_min, lo\_max)

    \EndFor

    \Statex

    \State Si hubo impacto (distancia menor a la máxima), marcar la celda:

    \If{distancia\_impacto < distancia\_maxima}

        \State celda\_impacto $\gets$
        celdas\_rayo[última]

        \State reticula[celda\_impacto] $\gets$
        saturar(

        \Statex \hspace{2.2cm}
        reticula[celda\_impacto] + incremento\_ocup,

        \Statex \hspace{2.2cm}
        lo\_min, lo\_max)

    \EndIf

\EndFor

\Statex
\State \Return reticula

\end{algorithmic}
\vspace{-0.9cm}
\rule{\textwidth}{0.4pt}

% -
\section{Pseudocódigo integrado de los módulos de representación}
\label{apend:representa}
% -

El siguiente pseudocódigo describe el ciclo completo de operación de los módulos de
representación, integrando SLAM y Ray Casting en un bucle continuo que corre mientras
el sistema está activo.

\begin{center}
\captionof{algorithm}{Pseudocódigo de representación}
\label{alg:representacion}
\end{center}
\vspace{-1.8cm}
\rule{\textwidth}{0.4pt}
\vspace{-0.9cm}

\begin{algorithmic}[1]

\Statex \textbf{INICIALIZACIÓN:}

\State pose $\gets$ [0, 0, 0]
\Comment{posición y orientación inicial}

\State covarianza $\gets$
matriz\_identidad * varianza\_inicial

\State reticula $\gets$
inicializar\_con\_log\_odds(0)

\Comment{todo desconocido}

\Statex

\While{sistema\_activo}

    \Statex \hspace{0.5cm}
    // Adquisición de datos del ciclo actual

    \State odometria $\gets$
    leer\_odometria()

    \Comment{(delta\_d, delta\_theta)}

    \State mediciones $\gets$
    leer\_sensores()

    \Comment{distancias por ángulo}

    \Statex

    \Statex \hspace{0.5cm}
    // Predicción de pose

    \State pose\_pred $\gets$
    aplicar\_modelo\_movimiento(pose, odometria)

    \State F $\gets$
    jacobiano\_movimiento(pose, odometria)

    \State covarianza\_pred $\gets$
    F * covarianza * F\textsuperscript{T} + Q

    \Statex

    \Statex \hspace{0.5cm}
    // Corrección con mediciones

    \State medicion\_esperada $\gets$
    modelo\_observacion(pose\_pred, reticula)

    \State innovacion $\gets$
    mediciones - medicion\_esperada

    \State H $\gets$
    jacobiano\_observacion(pose\_pred)

    \State S $\gets$
    H * covarianza\_pred * H\textsuperscript{T} + R

    \State K $\gets$
    covarianza\_pred * H\textsuperscript{T} * inversa(S)

    \State pose $\gets$
    pose\_pred + K * innovacion

    \State covarianza $\gets$
    (I - K * H) * covarianza\_pred

    \Statex

    \Statex \hspace{0.5cm}
    // Actualización del mapa con Ray Casting

    \For{cada rayo i}

        \State celdas $\gets$
        bresenham(pose, angulo[i], mediciones[i], resolucion)

        \For{c EN celdas excepto la última}

            \State reticula[c] $\gets$
            saturar(reticula[c] + lo\_libre, lo\_min, lo\_max)

        \EndFor

        \If{mediciones[i] < distancia\_maxima}

            \State reticula[celdas[última]] $\gets$
            saturar(

            \Statex \hspace{2.2cm}
            reticula[celdas[última]] + lo\_ocupado,

            \Statex \hspace{2.2cm}
            lo\_min, lo\_max)

        \EndIf

    \EndFor

    \Statex

    \Statex \hspace{0.5cm}
    // Publicar resultados para los demás módulos

    \State publicar(pose)

    \State publicar(reticula)

\EndWhile
\end{algorithmic}
\vspace{-0.9cm}
\rule{\textwidth}{0.4pt}

\section{Pseudocódigos Voz}

\subsubsection{Algoritmo 1: Detección de actividad de voz (VAD)}
\label{apend:vad}

La detección de actividad de voz separa los segmentos de audio que contienen voz de
los periodos de silencio o ruido de fondo, evitando procesar señal irrelevante.

\begin{center}
\captionof{algorithm}{Pseudocódigo VAD}
\label{alg:vad}
\end{center}
\vspace{-1.8cm}
\rule{\textwidth}{0.4pt}
\vspace{-0.9cm}

\begin{algorithmic}[1]

\Statex \textbf{ENTRADA:}
\Statex \hspace{0.5cm} flujo\_audio \Comment{flujo continuo de muestras de audio}

\Statex
\Statex \textbf{SALIDA:}
\Statex \hspace{0.5cm} segmentos\_voz \Comment{lista de fragmentos de audio con voz activa}

\Statex
\Statex \textbf{PROCESO}
\State Inicializar:
\State \hspace{0.5cm} buffer\_acumulado   $\gets$ vacío
\State \hspace{0.5cm} segmentos\_voz      $\gets$ []
\State \hspace{0.5cm} en\_voz             $\gets$ falso
\State \hspace{0.5cm} silencio\_acumulado $\gets$ 0 ms

\Statex
\For{cada bloque de 512 muestras EN flujo\_audio}

    \State Calcular energía del bloque
    \State energia $\gets$ promedio(cuadrado(muestras))

    \Statex
    \If{energia $>$ umbral\_inicio Y NO en\_voz}
        \State en\_voz             $\gets$ verdadero
        \State inicio\_segmento    $\gets$ tiempo\_actual()
        \State buffer\_acumulado   $\gets$ [bloque]
        \State silencio\_acumulado $\gets$ 0

    \Statex
    \ElsIf{en\_voz}
        \State agregar bloque a buffer\_acumulado

        \If{energia $<$ umbral\_fin}

            \State silencio\_acumulado $\gets$ silencio\_acumulado + duracion(bloque)

            \If{silencio\_acumulado $>$ pausa\_maxima (700 ms)}

                \State agregar buffer\_acumulado a segmentos\_voz
                \State en\_voz             $\gets$ falso
                \State silencio\_acumulado $\gets$ 0
                \State buffer\_acumulado   $\gets$ vacío

            \EndIf

        \Else

            \State silencio\_acumulado $\gets$ 0 \Comment{la voz continuó, reinciar pausa}

        \EndIf
\Statex
    \EndIf
\Statex
\EndFor
\Statex
\State \Return segmentos\_voz

\end{algorithmic}
\vspace{-0.9cm}
\rule{\textwidth}{0.4pt}

\subsubsection{Algoritmo 2: Reducción de ruido (Filtro de Wiener)}
\label{apend:wiener}
Sobre cada segmento de voz detectado se aplica un filtro espectral que estima el ruido
de fondo y lo elimina, mejorando la calidad de la señal antes de transcribirla.

Sea $X(f)$ el espectro del audio capturado, $N(f)$ el espectro estimado del ruido de
fondo y $G(f)$ la ganancia de Wiener aplicada a cada frecuencia $f$:
\[
G(f) = \frac{|X(f)|^2 - |N(f)|^2}{|X(f)|^2}
\]
La señal limpia se obtiene como:
\[
Y(f) = G(f) \cdot X(f)
\]
y se recupera mediante la transformada inversa.

\begin{center}
\captionof{algorithm}{Pseudocódigo reducción de ruido}
\label{alg:rudio}
\end{center}
\vspace{-1.8cm}
\rule{\textwidth}{0.4pt}
\vspace{-0.9cm}

\begin{algorithmic}[1]

\Statex \textbf{ENTRADA:}
\Statex segmento\_audio \Comment{fragmento de audio con voz}

\Statex
\Statex \textbf{SALIDA:}
\Statex audio\_limpio \Comment{audio con ruido reducido}

\Statex
\Statex \textbf{PROCESO}
\State Calcular el espectro de frecuencias del segmento
\State espectro $\gets$ transformada\_fourier(segmento\_audio)

\Statex
\State Estimar el espectro del ruido de fondo
\Comment{usando los primeros 100 ms sin voz activa}
\State ruido\_estimado $\gets$ estimar\_ruido\_fondo(espectro)

\Statex
\For{cada frecuencia f EN espectro}

    \State potencia\_señal $\gets$ $|espectro[f]|^2$
    \State potencia\_ruido $\gets$ $|ruido\_estimado[f]|^2$

    \Statex \hspace{0.5cm} // Calcular ganancia de Wiener (mínimo 0 para evitar amplificación)

    \State ganancia[f] $\gets$ maximo(0,
    \Statex \hspace{2.2cm} (potencia\_señal - potencia\_ruido) / potencia\_señal)

    \State espectro\_filtrado[f] $\gets$ ganancia[f] * espectro[f]

\EndFor

\Statex
\State Reconstruir señal en el dominio del tiempo
\State audio\_limpio $\gets$ transformada\_fourier\_inversa(espectro\_filtrado)

\Statex
\State \Return audio\_limpio

\end{algorithmic}
\vspace{-0.9cm}
\rule{\textwidth}{0.4pt}

OpenCV construye internamente una máscara gaussiana \(5 \times 5\) utilizando la siguiente ecuación:

\[
G(x,y)=\frac{1}{2\pi\sigma^2}e^{-\frac{x^2+y^2}{2\sigma^2}}
\]

donde:

\begin{itemize}
\item \(\sigma = 0.5\)
\item \(x,y\) representan las posiciones relativas respecto al centro del kernel
\end{itemize}

Para un kernel \(5 \times 5\) con \(\sigma = 0.5\), la distribución de pesos aproximada queda de la siguiente manera:

\[
\frac{1}{273}
\begin{bmatrix}
1 & 4 & 7 & 4 & 1 \\
4 & 16 & 26 & 16 & 4 \\
7 & 26 & 41 & 26 & 7 \\
4 & 16 & 26 & 16 & 4 \\
1 & 4 & 7 & 4 & 1
\end{bmatrix}
\]

\[
\begin{bmatrix}
0.0037 & 0.0147 & 0.0256 & 0.0147 & 0.0037 \\
0.0147 & 0.0586 & 0.0952 & 0.0586 & 0.0147 \\
0.0256 & 0.0952 & 0.1502 & 0.0952 & 0.0256 \\
0.0147 & 0.0586 & 0.0952 & 0.0586 & 0.0147 \\
0.0037 & 0.0147 & 0.0256 & 0.0147 & 0.0037
\end{bmatrix}
\]
